% Table 3: raw cross-tabulation and the cue-to-level mapping
% GENERATED by src/build_tables.py from data/numbers.json.
% Do not edit by hand; rerun the generator.
\begin{table}[pos=htbp]
\tableserif
\centering
\begin{threeparttable}
\hlcaption{HlM}{Coded severity against the narrative-implied injury indicator.}
\label{tab:audit}
\small
\renewcommand{\arraystretch}{1.15}
\begin{tabularx}{\textwidth}{l XXXX r}
\toprule
\rowcolor{PanelBg}
\multicolumn{6}{l}{\textbf{Panel A. Observed counts, with no model and no priors}} \\
\midrule
\headerrow
\hcell{KABCO} & \hcell{No cue} & \hcell{O-consistent} & \hcell{BC-consistent} & \hcell{KA-consistent} & \hcell{Total} \\
\midrule
O & \cellcolor{DiagGreen}50 (64.9) & 21 (27.3) & 6 (7.8) & 0 (0.0) & 77 \\
\altrow
C & \cellcolor{DiagGreen}69 (67.0) & 7 (6.8) & 26 (25.2) & 1 (1.0) & 103 \\
B & \cellcolor{DiagGreen}185 (71.2) & 16 (6.2) & 51 (19.6) & 8 (3.1) & 260 \\
\altrow
A & \cellcolor{DiagGreen}43 (60.6) & 0 (0.0) & 17 (23.9) & 11 (15.5) & 71 \\
K & 2 (18.2) & 0 (0.0) & 0 (0.0) & \cellcolor{DiagGreen}9 (81.8) & 11 \\
\midrule
\textbf{Total} & \textbf{349} & \textbf{44} & \textbf{100} & \textbf{29} & \textbf{522} \\
\bottomrule
\end{tabularx}

\vspace{6pt}
\begin{tabularx}{\textwidth}{c l X}
\toprule
\rowcolor{PanelBg}
\multicolumn{3}{l}{\textbf{Panel B. The mapping, applied in this order}} \\
\midrule
\headerrow
\hcell{Order} & \hcell{Level} & \hcell{Condition} \\
\midrule
1 & KA-consistent & fatality, fracture, or unconsciousness cue present \\
\altrow
2 & O-consistent & explicit no-injury cue present and no transport cue \\
3 & BC-consistent & transport cue present, no severe cue \\
\altrow
4 & no cue & none of the above \\
\bottomrule
\end{tabularx}
\begin{tablenotes}[flushleft]\footnotesize
\item[] \textit{Note:} Panel A gives counts with row percentages in parentheses, and
shaded cells hold more than 55\% of their row. Panel B is applied
in order, so a narrative with cues that point in two directions resolves the same way every
time. 15 narratives carried conflicting cues and were resolved by
this order.
\end{tablenotes}
\end{threeparttable}
\end{table}
